\documentclass[11pt,a4paper]{article}
\usepackage[T1]{fontenc}
\usepackage[utf8]{inputenc}
\usepackage[french]{babel}
\usepackage{lmodern}
\usepackage[a4paper,margin=19mm]{geometry}
\usepackage{microtype}
\usepackage{amsmath,amssymb}
\usepackage{enumitem}
\usepackage{longtable,tabularx}
\usepackage{xcolor}
\usepackage[hidelinks,breaklinks]{hyperref}
\definecolor{navy}{HTML}{18436B}
\definecolor{slate}{HTML}{39495B}
\definecolor{soft}{HTML}{E8F0F7}
\setlist[itemize]{leftmargin=*,itemsep=3pt,topsep=3pt}
\setlength{\parskip}{5pt}
\setlength{\parindent}{0pt}
\hypersetup{pdftitle={Dossier de references en mathematiques},pdfauthor={Dossier de travail, syntheses originales}}
\newcommand{\etat}[1]{\textsf{\textcolor{navy}{\textbf{[#1]}}}}
\newcommand{\refsource}[1]{\textsf{\textcolor{slate}{[#1]}}}
\newcommand{\sectionref}[2]{\section{#1}\textit{Références : #2.}\par}
\newcommand{\avertissement}[1]{\par\medskip\noindent\fcolorbox{navy}{soft}{\parbox{0.92\linewidth}{\textbf{Attention :} #1}}\par\medskip}
\newcommand{\entreesource}[5]{\par\noindent\textbf{\refsource{#1} #2}\\[-2pt]\textit{Statut : #3.} #4\\[-2pt]\href{#5}{\textcolor{navy}{\underline{Consulter le PDF officiel}}}\par\smallskip}
\newcommand{\entreelivre}[4]{\par\noindent\textbf{\refsource{#1} #2}\quad\textit{#3}\\[-2pt]#4\par\smallskip}
\renewcommand{\contentsname}{Sommaire}
\begin{document}
\begin{center}\LARGE\bfseries Dossier de références mathématiques --- 4e\end{center}
\begin{center}\small État des références : novembre 2026 --- Année scolaire 2026--2027\end{center}
\avertissement{Ce document est une synthèse de travail et un index de PDF, pas une reproduction des PDF. Les prescriptions viennent exclusivement du programme applicable. Les ressources pédagogiques anciennes et les propositions pour ChatMD ne créent aucune nouvelle exigence.}
\textbf{Programme de référence :} Programme de mathématiques du cycle 4 (2020). En vigueur en 4e en 2026-2027 ; le programme 2026 s’appliquera en 4e en 2027-2028.
\tableofcontents\newpage
\section{Comment lire ce dossier}
Ce dossier organise les sources \emph{sans en reproduire intégralement le contenu}. Chaque fiche synthétise un point de vigilance pour la préparation des cours ou pour un chatbot pédagogique. Les conseils opérationnels ne doivent pas être pris pour des citations littérales.\par
\textbf{Hiérarchie :} (1) programme en vigueur, (2) exemples/attendus de son année, (3) ressources didactiques, (4) propositions d’usage pour ChatMD. La 4e source, \emph{Des maths ensemble et pour chacun}, est complémentaire et peut refléter un ancien programme.\par
\section{Bibliothèque des PDF propres au niveau}
\entreesource{E4-01}{Programme de mathématiques du cycle 4 (2020)}{reglementaire}{Programme encore applicable en 4e en 2026-2027 ; programme 2026 seulement à partir de 2027-2028.}{https://eduscol.education.fr/document/621/download}
\entreesource{E4-02}{Attendus de fin d’année de mathématiques en 4e}{accompagnement actuel}{Niveau de maîtrise attendu, avec exemples de réussite.}{https://eduscol.education.gouv.fr/sites/default/files/document/16-maths-4e-attendus-eduscol1114746pdf-74682.pdf}
\entreesource{E4-03}{Repères annuels de progression en mathématiques au cycle 4}{accompagnement actuel}{Progressivité associée au programme 2020 pour les 4e et 3e.}{https://eduscol.education.gouv.fr/sites/default/files/document/26-maths-c4-reperes-eduscol1114756pdf-74694.pdf}
\entreesource{E4-04}{Utiliser le calcul littéral (ressource cycle 4)}{didactique ancien}{Expression procédurale et structurelle, variable, justification et équations.}{https://eduscol.education.gouv.fr/sites/default/files/document/ra16c4mathnombrescalculcalcullitteraldocmaitre548358pdf-77631.pdf}
\entreesource{E4-05}{Géométrie plane (ressource cycle 4)}{didactique ancien}{Propriétés, recherche de preuves et raisonnement géométrique.}{https://eduscol.education.gouv.fr/sites/default/files/document/ra16c4mathgeoplanedocmaitre574312pdf-77673.pdf}
\entreesource{E4-06}{Traitement des données (ressource cycle 4)}{didactique ancien}{Lecture, représentation, interprétation et comparaison de données.}{https://eduscol.education.gouv.fr/sites/default/files/document/ra16c4mathdocmaitre564036pdf-77637.pdf}
\sectionref{Référence réglementaire et limite temporelle}{E4-01, E4-02, E4-03}
\begin{itemize}
\item En 2026-2027, le texte de référence reste le programme de mathématiques du cycle 4 de 2020, complété par les attendus de fin d’année de 4e et les repères de progression applicables.
\item Le programme de cycle 4 publié le 5 mars 2026 entrera en vigueur en 4e en 2027-2028 : ne pas anticiper ses nouvelles exigences dans les évaluations de 2026-2027.
\item Pour toute notion sensible, toujours distinguer exigence de fin d’année, procédé d’enseignement choisi et approfondissement possible.
\end{itemize}
\sectionref{Pythagore et raisonnement géométrique}{E4-01, E4-02, E4-05, T-04}
\begin{itemize}
\item Le théorème de Pythagore ne s’applique qu’à un triangle rectangle ; identifier l’hypoténuse comme le côté opposé à l’angle droit.
\item Si l’énoncé donne seulement AB=6 cm et AC=8 cm dans ABC, la présence d’un angle droit n’est pas établie. Il faut demander si la figure ou l’énoncé l’indique ; ne pas supposer « rectangle en A ».
\item Dire « le triangle n’est pas rectangle en A » ne signifie pas « le triangle n’est pas rectangle » : il peut éventuellement être rectangle en B ou en C.
\item Le calcul d’un côté peut demander une addition ou une soustraction selon la position de l’inconnue. Distinguer choix de méthode, application correcte, calcul et rédaction.
\end{itemize}
\sectionref{Calcul littéral : égalité, structure, développement}{E4-01, E4-02, E4-04, T-05}
\begin{itemize}
\item Une expression a une structure algébrique : 3(x+5) est un produit d’un nombre par une somme ; 3x+5 est une somme de deux termes.
\item La distributivité s’applique pour développer le premier exemple, pas pour « développer » le second, déjà écrit sous forme de somme. Ne jamais inventer une parenthèse absente.
\item L’égalité entre deux expressions doit conserver la valeur pour toutes les valeurs pertinentes de la lettre. Dans une chaîne, ne jamais écrire un signe égal entre deux étapes non égales.
\item La lettre peut être nombre variable ou inconnue. Vérifier le rôle qu’elle prend dans la situation avant de choisir la procédure.
\end{itemize}
\sectionref{Résolution d’équations et statut de l’égalité}{E4-01, E4-02, E4-04}
\begin{itemize}
\item L’égalité est une relation entre deux expressions ; une transformation doit préserver l’ensemble des solutions. Pour 3x+5=17, retirer 5 aux deux membres est justifié ; écrire 3x=22 est incorrect.
\item Aider l’élève à identifier l’étape erronée et lui demander de la corriger, sans effectuer toutes les étapes restantes lorsqu’il demande seulement une vérification.
\item Faire distinguer recherche par essais, raisonnement et rédaction mathématique finale.
\end{itemize}
\sectionref{Fractions, relatifs, puissances, racines carrées}{E4-01, E4-02, E4-03, T-05}
\begin{itemize}
\item Conserver la distinction entre valeur exacte et valeur approchée : « = » n’a pas le même sens que «  $\approx$  ». Pour la racine carrée, contrôler le sens et l’encadrement avant une approximation.
\item Pour les calculs sur les relatifs, vérifier d’abord les opérations et les signes en jeu ; les parenthèses peuvent éviter toute ambiguïté.
\item Les automatismes sont importants, mais la procédure exigible doit être recoupée avec le programme 2020 et les attendus 4e, pas avec un manuel ancien.
\end{itemize}
\sectionref{Statistiques et traitement des données}{E4-01, E4-02, E4-06}
\begin{itemize}
\item Distinguer lire une donnée, calculer un indicateur et interpréter son sens dans la situation.
\item Pour une moyenne ou une médiane, préciser la série étudiée, l’unité et l’objet de la question. Éviter de présenter une mesure isolée comme une interprétation suffisante.
\item La représentation et les nombres doivent être comparés dans leur contexte ; le jugement sur une affirmation repose sur les données disponibles.
\end{itemize}
\sectionref{Rédaction, dialogue et autocontrôle}{T-01 à T-10, E4-05}
\begin{itemize}
\item Distinguer : données de l’énoncé ; hypothèses ; propriété applicable ; opérations ; réponse dans son contexte.
\item Préférer une question de relance décisive à une cascade de sous-questions qui effectuent le raisonnement à la place de l’élève.
\item Faire vérifier la cohérence d’un résultat, son unité et sa plausibilité. La trace rédigée n’est pas forcément identique aux essais de recherche.
\end{itemize}
\section{Références communes : dix documents transversaux}
\entreesource{T-01}{Chercher}{competence, didactique 2016}{Extraire des informations ; essayer, questionner et tester des pistes.}{https://eduscol.education.gouv.fr/sites/default/files/document/ra16c4mathchercher552125pdf-77583.pdf}

\entreesource{T-02}{Modéliser}{competence, didactique 2016}{Traduire une situation, examiner les hypothèses et la pertinence du modèle.}{https://eduscol.education.gouv.fr/sites/default/files/document/ra16c4mathmodelisernd566177pdf-77586.pdf}

\entreesource{T-03}{Représenter}{competence, didactique 2016}{Choisir et changer des registres de représentation.}{https://eduscol.education.gouv.fr/sites/default/files/document/ra16c4mathrepresenternd566207pdf-77589.pdf}

\entreesource{T-04}{Raisonner}{competence, didactique 2016}{Émettre, argumenter, justifier, vérifier des hypothèses.}{https://eduscol.education.gouv.fr/sites/default/files/document/ra16c4mathraisonner547836pdf-77592.pdf}

\entreesource{T-05}{Calculer}{competence, didactique 2016}{Conduire des calculs et contrôler les résultats.}{https://eduscol.education.gouv.fr/sites/default/files/document/ra16c4mathcompcalculer554370pdf-77595.pdf}

\entreesource{T-06}{Communiquer}{competence, didactique 2016}{Expliquer à l’oral et à l’écrit ; comprendre les formulations.}{https://eduscol.education.gouv.fr/sites/default/files/document/ra16c4mathcommecritoralpourmontage548548pdf-77598.pdf}

\entreesource{T-07}{Guide Résolution de problèmes mathématiques au collège}{guide transversal}{Résolution, erreurs, stratégies et interprétation des résultats.}{https://eduscol.education.gouv.fr/sites/default/files/document/guide-resolution-de-problemes-mathematiques-au-college-73914.pdf}

\entreesource{T-08}{Trace écrite de cours en mathématiques}{guide transversal}{Fonctions de la trace écrite, construction et organisation du cours.}{https://eduscol.education.gouv.fr/sites/default/files/document/trace-ecrite-de-cours-en-mathspdf-91731.pdf}

\entreesource{T-09}{Les pratiques orales dans l’enseignement des mathématiques}{guide transversal}{Oral de recherche, explication, échanges et relances.}{https://eduscol.education.gouv.fr/sites/default/files/document/les-pratiques-orales-dans-lenseignement-des-mathspdf-91728.pdf}

\entreesource{T-10}{Les automatismes}{guide transversal}{Développer, entretenir et rendre disponibles les connaissances de base.}{https://eduscol.education.gouv.fr/sites/default/files/document/automatismes-mathematiquespdf-91875.pdf}
\subsection{Les six compétences : repères de lecture pour ChatMD}
\begin{itemize}
\item \textbf{Chercher :} identifier des informations pertinentes, formuler une question, essayer une piste ; ne pas imposer trop tôt une méthode.
\item \textbf{Modéliser :} expliciter le modèle choisi et ses conditions d'emploi ; distinguer fait observé et hypothèse.
\item \textbf{Représenter :} choisir schéma, tableau, droite graduée, figure ou écriture ; relier les représentations.
\item \textbf{Raisonner :} justifier une affirmation à partir des données et des propriétés, ou produire un contre-exemple.
\item \textbf{Calculer :} effectuer une transformation valide puis contrôler signes, unités, ordre de grandeur et égalité.
\item \textbf{Communiquer :} rendre une démarche intelligible ; distinguer échanges, essais et rédaction mathématique stabilisée.
\end{itemize}
\textit{Ces repères sont une reformulation de travail issue des ressources transversales, et non une grille officielle d'évaluation.}

\section{Mode d’emploi dans Overleaf et ChatMD}
\begin{itemize}\item Ce fichier \texttt{.tex} est compilable de façon autonome sous pdfLaTeX.\item Les liens conduisent aux PDF officiels : ils ne téléchargent pas les PDF dans le projet Overleaf.\item Pour que ChatMD exploite ce contenu, intégrer le texte dans l’index RAG ; la simple présence du \texttt{.tex} dans Overleaf ne suffit pas.\item Revalider le calendrier réglementaire à chaque rentrée scolaire.\end{itemize}
\end{document}
